\begin{lemma}
\label{editorial-source-lemma-help-with-powers}
Let $p$ be a prime number. Let $n, m > 0$ be two integers. There exists
an integer $a$ such that
$(x + y)^{p^a}, p^a(x + y) \in \mathbf{Z}[x^{p^n}, p^nx, y^{p^m}, p^my]$.
\end{lemma}

\begin{proof}
Retain the original subring
$$
A=\mathbf Z[x^{p^n},p^n x,y^{p^m},p^m y]
\subseteq\mathbf Z[x,y].
$$
We calculate exactly which integer
coefficients can occur at each original
monomial $x^i y^j$. Write uniquely
$$
i=q p^n+r,\quad 0\leq r<p^n,
\qquad j=q' p^m+r',\quad 0\leq r'<p^m.
$$
In any product of the four displayed
generators contributing to this
monomial, the exponent on $p^n x$
is at least $r$, and the exponent
on $p^m y$ is at least $r'$.
Thus its coefficient is divisible
by $p^{nr+mr'}$. Conversely the
product
$$
(x^{p^n})^q(p^n x)^r
(y^{p^m})^{q'}(p^m y)^{r'}
=p^{nr+mr'}x^i y^j
$$
belongs to $A$. Since distinct formal
monomials are independent over
$\mathbf Z$, a polynomial belongs
to $A$ exactly when each of its
coefficients is divisible by this
corresponding integer. This proves
necessity as well as sufficiency
of the source divisibility test.

For a nonnegative integer $a$,
$p^a(x+y)$ belongs to $A$ exactly
when $a\geq b=\max(n,m)$: the
coefficient tests for $x$ and $y$
are respectively divisibility by
$p^n$ and $p^m$. In that case its
full expression in the generators is
$$
p^a(x+y)=p^{a-n}(p^n x)+p^{a-m}(p^m y).
$$
Now keep $a\geq b$ and $N=p^a$.
For $0<i<N$, the original binomial
coefficient has the exact valuation
$$
v_p\binom Ni=a-v_p(i).
$$
Indeed
$\binom Ni=(N/i)\binom{N-1}{i-1}$,
and the second factor is not divisible
by $p$. To verify the latter assertion,
reduce the full polynomial identity
$(1+T)^N=1+T^N$ modulo $p$ and divide
by $1+T$ in $\mathbf F_p[T]$.
It gives
$$
(1+T)^{N-1}=\sum_{h=0}^{N-1}(-T)^h
\quad\hbox{in }\mathbf F_p[T].
$$
Multiplication by $1+T$ verifies this
identity: the last sign is positive
for odd $p$, and signs coincide for
$p=2$. Every coefficient is nonzero,
giving the valuation assertion with
the original integer coefficient
still retained.

The two endpoint terms $i=0,N$
already belong to $A$. For an
intermediate $i$, put $t=i\bmod p^b$
with $0\leq t<p^b$. If $t=0$, both
remainders $r,r'$ are zero. If $t>0$,
then $v_p(i)=v_p(t)$ and
$$
r=t\bmod p^n,\qquad
r'=(-t)\bmod p^m,
$$
where both residues are chosen in
their original nonnegative ranges.
Every $1\leq t<p^b$ occurs among
$1\leq i<N$. Consequently the least
possible $a$ for both assertions
is exactly
$$
\max\left(b,
\max_{1\leq t<p^b}
\left(v_p(t)+n(t\bmod p^n)+m((-t)\bmod p^m)\right)\right).
$$
This is an exact test on the
original integer coefficients,
not merely a sufficient estimate.

It has the following explicit value:
$$
a_{min}=
\begin{cases}
n p^n+n-1,&n=m,\\
b(p^b-1)+c,&n\neq m,\quad
b=\max(n,m),\ c=\min(n,m).
\end{cases}
$$
For $n=m$, each nonzero residue $t$
has $r+r'=p^n$, and its valuation
is at most $n-1$, with equality at
$t=p^{n-1}$. This gives the first
value, which is at least $b$.
For $n<m$, write $t=q p^n+r$.
If $r>0$, its full weight is
$$
m p^m-m q p^n-(m-n)r+v_p(r).
$$
For every positive $r$ one has
$v_p(r)\leq r-1$: if $v=v_p(r)$,
then $r\geq p^v\geq v+1$, the
last inequality following by
induction on $v$. Since $m-n\geq1$,
this gives
$v_p(r)\leq(m-n)(r-1)$.
The weight is therefore at most
$m p^m-(m-n)$, and equality is
attained at the original $t=1$.
If $r=0$ and $t>0$, its weight is
$m(p^m-t)+v_p(t)$, at most
$m p^m-(m-1)t-1$. Since $t\geq1$,
this is at most $m p^m-m$, and
hence at most $m p^m-(m-n)$.
Thus the maximum is exactly
$m(p^m-1)+n$. For $m<n$, the
automorphism of $\mathbf Z[x,y]$
exchanging the original $x$ and
$y$ takes $A$ to the corresponding
subring with $n,m$ exchanged, and
fixes both $(x+y)^{p^a}$ and
$p^a(x+y)$. Applying this actual
inverse-to-itself map proves the
other unequal case. These values
are all at least $b$.

For every $a\geq a_{min}$, the full
desired expression is therefore
$$
(x+y)^{p^a}=
\sum_{\substack{i,j\geq0\\i+j=p^a}}
\frac{\binom{p^a}{i,j}}{p^{nr+mr'}}
(x^{p^n})^q(p^n x)^r
(y^{p^m})^{q'}(p^m y)^{r'}.
$$
Each displayed quotient is an
integer by the proved divisibility;
no power of $p$ is inverted in the
ring. All original binomial
coefficients and all four generator
factors are retained. If $b\leq a<a_{min}$,
the maximizing residue just described
gives an actual monomial whose
coefficient fails the necessary
divisibility. If $a<b$, the linear
expression already fails. This
proves minimality.

In particular the source choice
$a_0=n p^n+m p^m+n+m$ works.
For $n=m$, its difference from
$a_{min}$ is $n p^n+n+1$.
For $n\neq m$, with the same
$b,c$, the difference is $c p^c+2b$.
Both are positive. Thus the original
sufficient choice is preserved and
verified, alongside the sharp bound.
\end{proof}